\documentclass[10pt,a4paper]{amsart}
\usepackage{amsmath,amssymb,amsthm,graphicx}
\usepackage{verbatim}
\newenvironment{reference}{\comment}{\endcomment}
%\newenvironment{reference}{}{}
\newenvironment{slogan}{\comment}{\endcomment}
\newenvironment{history}{\comment}{\endcomment}

% For commutative diagrams we use Xy-pic
\usepackage[all]{xy}

% We use 2cell for 2-commutative diagrams.
\xyoption{2cell}
% Theorem environments.
%
\theoremstyle{plain}
\newtheorem{theorem}[subsection]{Theorem}
\newtheorem{proposition}[subsection]{Proposition}
\newtheorem{lemma}[subsection]{Lemma}

\theoremstyle{definition}
\newtheorem{definition}[subsection]{Definition}
\newtheorem{example}[subsection]{Example}
\newtheorem{exercise}[subsection]{Exercise}
\newtheorem{situation}[subsection]{Situation}

\theoremstyle{remark}
\newtheorem{remark}[subsection]{Remark}
\newtheorem{remarks}[subsection]{Remarks}

\numberwithin{equation}{subsection}

% Macros
%
\def\lim{\mathop{\mathrm{lim}}\nolimits}
\def\colim{\mathop{\mathrm{colim}}\nolimits}
\def\Spec{\mathop{\mathrm{Spec}}}
\def\Hom{\mathop{\mathrm{Hom}}\nolimits}
\def\Ext{\mathop{\mathrm{Ext}}\nolimits}
\def\SheafHom{\mathop{\mathcal{H}\!\mathit{om}}\nolimits}
\def\SheafExt{\mathop{\mathcal{E}\!\mathit{xt}}\nolimits}
\def\Sch{\mathit{Sch}}
\def\Mor{\mathop{\mathrm{Mor}}\nolimits}
\def\Ob{\mathop{\mathrm{Ob}}\nolimits}
\def\Sh{\mathop{\mathit{Sh}}\nolimits}
\def\NL{\mathop{N\!L}\nolimits}
\def\CH{\mathop{\mathrm{CH}}\nolimits}
\def\proetale{{pro\text{-}\acute{e}tale}}
\def\etale{{\acute{e}tale}}
\def\QCoh{\mathit{QCoh}}
\def\Ker{\mathop{\mathrm{Ker}}}
\def\Im{\mathop{\mathrm{Im}}}
\def\Coker{\mathop{\mathrm{Coker}}}
\def\Coim{\mathop{\mathrm{Coim}}}

% Boxtimes
%
\DeclareMathSymbol{\boxtimes}{\mathbin}{AMSa}{"02}

%
% Macros for moduli stacks/spaces
%
\def\QCohstack{\mathcal{QC}\!\mathit{oh}}
\def\Cohstack{\mathcal{C}\!\mathit{oh}}
\def\Spacesstack{\mathcal{S}\!\mathit{paces}}
\def\Quotfunctor{\mathrm{Quot}}
\def\Hilbfunctor{\mathrm{Hilb}}
\def\Curvesstack{\mathcal{C}\!\mathit{urves}}
\def\Polarizedstack{\mathcal{P}\!\mathit{olarized}}
\def\Complexesstack{\mathcal{C}\!\mathit{omplexes}}
% \Pic is the operator that assigns to X its picard group, usage \Pic(X)
% \Picardstack_{X/B} denotes the Picard stack of X over B
% \Picardfunctor_{X/B} denotes the Picard functor of X over B
\def\Pic{\mathop{\mathrm{Pic}}\nolimits}
\def\Picardstack{\mathcal{P}\!\mathit{ic}}
\def\Picardfunctor{\mathrm{Pic}}
\def\Deformationcategory{\mathcal{D}\!\mathit{ef}}
\usepackage{cleveref}
\crefname{originallemma}{original lemma}{original lemmas}
\Crefname{originallemma}{Original lemma}{Original lemmas}
\makeatletter
\newlabel{coherent-lemma-chow-Noetherian}{{0200}{source}{}{sourcepointer-1}{}}
\newlabel{coherent-lemma-chow-Noetherian@cref}{{[originallemma][0][]0200}{[1][0][]source}{}{}{}}
\makeatother
\crefname{originallemma}{original lemma}{original lemmas}
\Crefname{originallemma}{Original lemma}{Original lemmas}
\makeatletter
\newlabel{divisors-lemma-blowing-up-projective}{{02NS}{source}{}{sourcepointer-2}{}}
\newlabel{divisors-lemma-blowing-up-projective@cref}{{[originallemma][0][]02NS}{[1][0][]source}{}{}{}}
\makeatother
\crefname{originalreference}{original reference}{original references}
\Crefname{originalreference}{Original reference}{Original references}
\makeatletter
\newlabel{lemma-contract-when-quasi-projective}{{0C2N}{source}{}{sourcepointer-3}{}}
\newlabel{lemma-contract-when-quasi-projective@cref}{{[originalreference][0][]0C2N}{[1][0][]source}{}{}{}}
\makeatother
\crefname{originalreference}{original reference}{original references}
\Crefname{originalreference}{Original reference}{Original references}
\makeatletter
\newlabel{lemma-contraction-unique}{{0C5K}{source}{}{sourcepointer-4}{}}
\newlabel{lemma-contraction-unique@cref}{{[originalreference][0][]0C5K}{[1][0][]source}{}{}{}}
\makeatother
\crefname{originalreference}{original reference}{original references}
\Crefname{originalreference}{Original reference}{Original references}
\makeatletter
\newlabel{lemma-extend-rational-map-blowing-up}{{0C5H}{source}{}{sourcepointer-5}{}}
\newlabel{lemma-extend-rational-map-blowing-up@cref}{{[originalreference][0][]0C5H}{[1][0][]source}{}{}{}}
\makeatother
\crefname{originallemma}{original lemma}{original lemmas}
\Crefname{originallemma}{Original lemma}{Original lemmas}
\makeatletter
\newlabel{morphisms-lemma-composition-projective}{{0C4P}{source}{}{sourcepointer-6}{}}
\newlabel{morphisms-lemma-composition-projective@cref}{{[originallemma][0][]0C4P}{[1][0][]source}{}{}{}}
\makeatother
\crefname{originallemma}{original lemma}{original lemmas}
\Crefname{originallemma}{Original lemma}{Original lemmas}
\makeatletter
\newlabel{morphisms-lemma-composition-quasi-projective}{{0C4M}{source}{}{sourcepointer-7}{}}
\newlabel{morphisms-lemma-composition-quasi-projective@cref}{{[originallemma][0][]0C4M}{[1][0][]source}{}{}{}}
\makeatother
\crefname{originallemma}{original lemma}{original lemmas}
\Crefname{originallemma}{Original lemma}{Original lemmas}
\makeatletter
\newlabel{morphisms-lemma-projective-permanence}{{0C4Q}{source}{}{sourcepointer-8}{}}
\newlabel{morphisms-lemma-projective-permanence@cref}{{[originallemma][0][]0C4Q}{[1][0][]source}{}{}{}}
\makeatother
\crefname{originallemma}{original lemma}{original lemmas}
\Crefname{originallemma}{Original lemma}{Original lemmas}
\makeatletter
\newlabel{morphisms-lemma-projective-quasi-projective}{{07RL}{source}{}{sourcepointer-9}{}}
\newlabel{morphisms-lemma-projective-quasi-projective@cref}{{[originallemma][0][]07RL}{[1][0][]source}{}{}{}}
\makeatother
\title{Quasi-projectivity of a regular two dimensional scheme over a Noetherian base}
\author{The Stacks Project Authors}
\date{stacks-project:resolve.tex:local-snapshot-20261012; GNU FDL 1.2 or later}
\begin{document}
\maketitle
\noindent\textit{Editorial note.} All delivered mathematical text is the author's own, copied from the pinned Stacks Project chapter file resolve.tex; no statement or proof step was rewritten, reordered or removed.\par\medskip
\noindent\textit{Editorial note.} The counted proof is the author's complete argument: the application of Chow's lemma, the extension of the rational map along blowups in closed points, and the descent of quasi-projectivity along the chain of blowups.\par\medskip
\noindent\textit{Editorial note.} The label lemma-regular-dim-2-quasi-projective names the counted statement, and the proof environment immediately following it in the source is delivered verbatim as its proof.\par\medskip
\noindent\textit{Editorial note.} The counted statement is carried by the author's own lemma environment in the chapter source.\par\medskip
\noindent\textit{Editorial note.} Pre-open gate: python operations/agent-authoring/tools/precard\_gate.py --source resolve.tex --label lemma-regular-dim-2-quasi-projective --start 5049 --end 5080 --local-id 28105 returned RESULT=OK (L28105 free, occupancy FREE, public interval clear). Pre-assembly tag self-check: canonical tag resolved and not present in reference/public-items.jsonl.\par\medskip
\noindent\textit{Editorial note.} The statement and proof are copied from the pinned chapter file; the cited inputs are kept as external pointers with their canonical Stacks tags.\par\medskip
\noindent\textit{Allowed standard prerequisites (editorial).}\begin{itemize}
\item Chow's lemma. Cohomology of Schemes, the Noetherian version of Chow's lemma. Produces a proper modification of the source which is H-quasi-projective over the base and an isomorphism over a dense open.
\item Extending a rational map along a sequence of blowups in closed points. Same chapter, the lemma producing a sequence of blowups in closed points lifting a rational map into a modification. Provides the chain of blowups through which the rational map to the H-quasi-projective modification is defined.
\item Permanence of projectivity and quasi-projectivity, and contraction of the last blowing up. Morphisms, the lemmas on the composition and permanence of projective and quasi-projective morphisms, together with the lemma on contracting the last blowup when the first blowup is quasi-projective and the uniqueness of contractions. Transfers quasi-projectivity from the last blowup back down the chain to the original scheme.
\end{itemize}\medskip
\ifdefined\hypertarget\hypertarget{sourcepointer-1}{}\fi
\noindent\textit{Source reference.} Original lemma 0200: sources/stacks-chow/tags.txt, source lines 2617--2617. Prerequisite pointer invoked at proof line(s) 5058; it names a standard fact of a sibling chapter that the counted argument uses, and every step of the counted proof built on it is retained verbatim in the delivered body.\par\medskip
\ifdefined\hypertarget\hypertarget{sourcepointer-2}{}\fi
\noindent\textit{Source reference.} Original lemma 02NS: sources/stacks-chow/tags.txt, source lines 3491--3491. Prerequisite pointer invoked at proof line(s) 5068; it names a standard fact of a sibling chapter that the counted argument uses, and every step of the counted proof built on it is retained verbatim in the delivered body.\par\medskip
\ifdefined\hypertarget\hypertarget{sourcepointer-3}{}\fi
\noindent\textit{Source reference.} Original reference 0C2N: sources/stacks-resolve/pinned-original-source.tex, source lines 4986--4986. Prerequisite pointer invoked at proof line(s) 5076; it names a standard fact of the same chapter that the counted argument uses, and every step of the counted proof built on it is retained verbatim in the delivered body.\par\medskip
\ifdefined\hypertarget\hypertarget{sourcepointer-4}{}\fi
\noindent\textit{Source reference.} Original reference 0C5K: sources/stacks-resolve/pinned-original-source.tex, source lines 4592--4592. Prerequisite pointer invoked at proof line(s) 5077; it names a standard fact of the same chapter that the counted argument uses, and every step of the counted proof built on it is retained verbatim in the delivered body.\par\medskip
\ifdefined\hypertarget\hypertarget{sourcepointer-5}{}\fi
\noindent\textit{Source reference.} Original reference 0C5H: sources/stacks-resolve/pinned-original-source.tex, source lines 1182--1182. Prerequisite pointer invoked at proof line(s) 5061; it names a standard fact of the same chapter that the counted argument uses, and every step of the counted proof built on it is retained verbatim in the delivered body.\par\medskip
\ifdefined\hypertarget\hypertarget{sourcepointer-6}{}\fi
\noindent\textit{Source reference.} Original lemma 0C4P: sources/stacks-chow/tags.txt, source lines 15234--15234. Prerequisite pointer invoked at proof line(s) 5069; it names a standard fact of a sibling chapter that the counted argument uses, and every step of the counted proof built on it is retained verbatim in the delivered body.\par\medskip
\ifdefined\hypertarget\hypertarget{sourcepointer-7}{}\fi
\noindent\textit{Source reference.} Original lemma 0C4M: sources/stacks-chow/tags.txt, source lines 15232--15232. Prerequisite pointer invoked at proof line(s) 5073; it names a standard fact of a sibling chapter that the counted argument uses, and every step of the counted proof built on it is retained verbatim in the delivered body.\par\medskip
\ifdefined\hypertarget\hypertarget{sourcepointer-8}{}\fi
\noindent\textit{Source reference.} Original lemma 0C4Q: sources/stacks-chow/tags.txt, source lines 15235--15235. Prerequisite pointer invoked at proof line(s) 5071; it names a standard fact of a sibling chapter that the counted argument uses, and every step of the counted proof built on it is retained verbatim in the delivered body.\par\medskip
\ifdefined\hypertarget\hypertarget{sourcepointer-9}{}\fi
\noindent\textit{Source reference.} Original lemma 07RL: sources/stacks-chow/tags.txt, source lines 9871--9871. Prerequisite pointer invoked at proof line(s) 5075; it names a standard fact of a sibling chapter that the counted argument uses, and every step of the counted proof built on it is retained verbatim in the delivered body.\par\medskip
\begin{lemma}
\label{lemma-regular-dim-2-quasi-projective}
Let $S$ be a Noetherian scheme. Let $f : X \to S$ be a
separated morphism of finite type with $X$ regular of dimension $2$.
Then $X$ is quasi-projective over $S$.
\end{lemma}

\begin{proof}
By Chow's lemma
(Cohomology of Schemes, Lemma \ref{coherent-lemma-chow-Noetherian})
there exists a proper morphism $\pi : X' \to X$ which is an isomorphism
over a dense open $U \subset X$ such that $X' \to S$ is H-quasi-projective.
By Lemma \ref{lemma-extend-rational-map-blowing-up}
there exists a sequence of blowups in closed points
$$
X_n \to \ldots \to X_1 \to X_0 = X
$$
and an $S$-morphism $X_n \to X'$ extending the rational map $U \to X'$.
Observe that $X_n \to X$ is projective by
Divisors, Lemma \ref{divisors-lemma-blowing-up-projective} and
Morphisms, Lemma \ref{morphisms-lemma-composition-projective}.
This implies that $X_n \to X'$ is projective by
Morphisms, Lemma \ref{morphisms-lemma-projective-permanence}.
Hence $X_n \to S$ is quasi-projective by
Morphisms, Lemma \ref{morphisms-lemma-composition-quasi-projective}
(and the fact that a projective morphism is quasi-projective, see
Morphisms, Lemma \ref{morphisms-lemma-projective-quasi-projective}).
By Lemma \ref{lemma-contract-when-quasi-projective}
(and uniqueness of contractions Lemma \ref{lemma-contraction-unique})
we conclude  that $X_{n - 1}, \ldots, X_0 = X$ are quasi-projective over $S$
as desired.
\end{proof}

\end{document}
